\import{set/powerset.tex}
\import{set/fixpoint.tex}
\import{function.tex}

\section{Equinumerosity}


\begin{definition}\label{equinum}
    $A$ is equinumerous with $B$ iff there exists a bijection from $A$ to $B$.
\end{definition}

\begin{abbreviation}\label{approx}
    $A\approx B$ iff $A$ is equinumerous with $B$.
\end{abbreviation}

\begin{proposition}\label{equinum_refl}
    $A\approx A$.
\end{proposition}
\begin{proof}
    $\identity{A}$ is a bijection from $A$ to $A$ by \obv{id_elem_bijections}.
    Follows by \obv{equinum}.
\end{proof}

\begin{proposition}\label{equinum_sym}
    Suppose $A\approx B$. Then $B\approx A$.
\end{proposition}
\begin{proof}
    Take $f$ such that $f$ is a bijection from $A$ to $B$ by \obv{equinum}.
    Then $\converse{f}$ is a bijection from $B$ to $A$ by \obv{bijection_converse_is_bijection}.
    Follows by \obv{equinum}.
\end{proof}

\begin{proposition}\label{equinum_inhabited}
    Suppose $A\approx B$ and $B$ is inhabited.
    Then $A$ is inhabited.
\end{proposition}
\begin{proof}
    Take $f$ such that $f\in\bijections{A}{B}$ by \obv{equinum}.
    Hence $f\in\surj{A}{B}$ by \obv{bijections}.
    Take $b$ such that $b\in B$.
    Take $a\in A$ such that $f(a)=b$ by \obv{surj}.
    Follows by \auto.
\end{proof}

\begin{proposition}\label{equinum_tran}
    Suppose $A\approx B\approx C$. Then $A\approx C$.
\end{proposition}
\begin{proof}
    Take $f$ such that $f$ is a bijection from $A$ to $B$ by \obv{equinum}.
    Take $g$ such that $g$ is a bijection from $B$ to $C$ by \obv{equinum}.
    Then $g\circ f$ is a bijection from $A$ to $C$ by \obv{bijection_circ}.
    Follows by \obv{equinum}.
\end{proof}



\begin{theorem}[Cantor--Schröder--Bernstein]\label{cantorschroederbernstein}
    Let $f$ be an injective function from $A$ to $B$.
    Let $g$ be an injective function from $B$ to $A$.
    Then $A\approx B$.
\end{theorem}
\begin{proof}
    $f\in\funs{A}{B}$ and $g\in\funs{B}{A}$.
    Thus $f$ and $g$ are functions by \obv{funs_is_function}.
    $\dom{f}=A$ and $\dom{g}=B$ by \obv{funs}.
    $f$ and $g$ are injective.
    $f\in\inj{A}{B}$ and $g\in\inj{B}{A}$
        by \auto{inj,injective_function}.

    For all $X\in\pow{A}$ we have
        $A\setminus \img{g}{B\setminus\img{f}{X}}\subseteq A$
        by \obv{setminus_subseteq}.
    For all $X\in\pow{A}$ we have
        $A\setminus \img{g}{B\setminus\img{f}{X}}\in\pow{A}$
        by \obv{powerset_intro}.
    Let $h(X) = A\setminus \img{g}{B\setminus\img{f}{X}}$ for $X\in\pow{A}$.
    For all $X\in\pow{A}$ we have
        $(X,A\setminus \img{g}{B\setminus\img{f}{X}})\in h$.
    $h$ is a relation by \auto{relation}.
    $h$ is right-unique by \auto{rightunique,pair_eq_iff}.
    For all $X\in\pow{A}$ we have
        $h(X) = A\setminus \img{g}{B\setminus\img{f}{X}}$
        by \obv{function_apply_intro}.
    For all $X\in\pow{A}$ we have $h(X)\in\pow{A}$ by \auto.
    For all $X$ we have if $X\in\dom{h}$, then $X\in\pow{A}$
        by \auto{dom_iff,pair_eq_iff}.
    For all $X$ we have if $X\in\pow{A}$, then $X\in\dom{h}$
        by \auto{dom_iff}.
    Thus $\dom{h} = \pow{A}$ by set extensionality.
    $h$ is a function to $\pow{A}$.
    Thus $h$ is a function from $\pow{A}$ to $\pow{A}$ by \obv{funs_intro}.

    $h$ is \subseteq-preserving by \auto{subseteqpreserving,img_subseteq,subseteq_implies_setminus_supseteq}. % apply each each lemma twice (alternating).

    Take $X$ such that $X$ is a fixpoint of $h$ by \obv{knastertarski}.
    $X\in\dom{h}$ and $h(X)=X$ by \obv{fixpoint}.
    Thus $X\in\pow{A}$.
    Hence $X\subseteq A$ by \obv{powerset_elim}.
    \begin{align*}
        X
            &= h(X)
            \\
            &= A\setminus \img{g}{B\setminus\img{f}{X}}
    \end{align*}

    $\img{g}{B\setminus\img{f}{X}}\subseteq A$
        by \auto{img_subseteq_ran,funs_ran,subseteq}.

    Thus $\img{g}{B\setminus\img{f}{X}} = A\setminus X$ by \auto{double_relative_complement}.

    Let $p=\restrl{f}{X}$.
    $p\in\inj{X}{B}$ by \auto{inj_restrl}.
    $\ran{p}=\img{f}{X}$ by \auto{restrl_ran}.

    $A\setminus X\subseteq\ran{g}$
        by \auto{img_subseteq_ran,subseteq}.
    $\converse{g}\in\inj{\ran{g}}{B}$
        by \obv{converse_inj_ran_dom}.
    Let $q=\restrl{\converse{g}}{A\setminus X}$.
    $q\in\inj{A\setminus X}{B}$ by \auto{inj_restrl}.

    Show $\ran{q}=B\setminus\img{f}{X}$.
    \begin{subproof}
        Show every element of $\ran{q}$ is an element of
            $B\setminus\img{f}{X}$.
        \begin{subproof}
            Fix $b\in\ran{q}$.
            Take $a$ such that $(a,b)\in q$ by \obv{ran_iff}.
            Then $a\in A\setminus X$ and $(a,b)\in\converse{g}$
                by \auto{restrl_iff}.
            Thus $(b,a)\in g$ by \obv{converse_iff}.
            Take $c\in B\setminus\img{f}{X}$ such that $(c,a)\in g$
                by \auto{img_iff}.
            $b=c$ by \obv{injective}.
            Thus $b\in B\setminus\img{f}{X}$.
        \end{subproof}
        Show every element of $B\setminus\img{f}{X}$ is an element of
            $\ran{q}$.
        \begin{subproof}
            Fix $b\in B\setminus\img{f}{X}$.
            $b\in B$ by \obv{setminus}.
            Thus $b\in\dom{g}$ by \obv{funs}.
            $(b,g(b))\in g$ by \obv{function_apply_elim}.
            Hence $g(b)\in\img{g}{B\setminus\img{f}{X}}$
                by \auto{img_iff}.
            Thus $g(b)\in A\setminus X$.
            $(g(b),b)\in\converse{g}$ by \obv{converse_iff}.
            Hence $(g(b),b)\in q$ by \auto{restrl_iff}.
            Thus $b\in\ran{q}$ by \auto{ran_iff}.
        \end{subproof}
        Follows by set extensionality.
    \end{subproof}

    $X\inter(A\setminus X)=\emptyset$ by \obv{setminus_disjoint}.
    $\ran{p}\inter\ran{q}=\emptyset$
        by \auto{setminus_disjoint}.
    Let $k=p\union q$.
    $k\in\inj{X\union(A\setminus X)}{B}$
        by \auto{union_inj_disjoint}.
    $X\union(A\setminus X)=A$ by \obv{setminus_partition}.
    Thus $k\in\inj{A}{B}$.

    $\img{f}{X}\subseteq B$
        by \auto{img_subseteq_ran,funs_ran,subseteq}.
    $\img{f}{X}\union(B\setminus\img{f}{X})=B$
        by \obv{setminus_partition}.
    $\ran{k}=\ran{p}\union\ran{q}$ by \auto{ran_union}.
    Thus $\ran{k}=B$.
    $k\in\funs{A}{B}$ by \obv{inj}.
    $k\in\surj{A}{B}$ by \obv{funs_surj_iff}.
    $k$ is a function by \obv{funs_is_function}.
    $\dom{k}=A$ by \obv{funs}.
    For all $x,y\in\dom{k}$ such that $k(x)=k(y)$ we have $x=y$
        by \obv{inj}.
    $k$ is injective by \obv{injective_function}.
    Hence $k\in\bijections{A}{B}$ by \obv{bijections}.

    Follows by \obv{equinum}.
\end{proof}
